%% =====================================================================
%%  Decomposing multi-hop questions into ordered atomic steps
%%  Progress review, September 2026 -- Jahid Ahmadi, Sabanci University
%%
%%  Beamer transcription of Slides/thesis-deck-v2.html  (35 slides).
%%  Compile with pdflatex.  Self-contained: paste into a blank Overleaf
%%  project as the only .tex file and hit Recompile.  No images, no
%%  \input, no bibliography.
%% =====================================================================
\documentclass[aspectratio=169,11pt]{beamer}

\usepackage[T1]{fontenc}
\usepackage[utf8]{inputenc}
\usepackage{microtype}
\usepackage{xcolor}
\usepackage{tikz}
\usepackage{booktabs}
\usepackage{array}
\usepackage{amsmath}

% ---------------------------------------------------------------- fonts
% If plex-sans is unavailable, comment the two lines below and uncomment this:
% \usepackage{lmodern}
\usepackage{plex-sans}
\usepackage{plex-mono}

% Math in Plex.  Beamer already sets math letters in the sans text font; this
% pins the digits, the decimal point and \mathrm to Plex Sans as well, so a
% number reads the same inside and outside a formula.  The T1 IBMPlexSans-TLF
% family ships with plex-sans and exists under pdflatex and XeTeX alike.
% Learnings (metric formulas, 2026-10-14):
%  - under XeTeX, fontspec moves math digits to the serif \rmdefault at
%    \begin{document}; re-declaring the digit math symbols (and \mathbf,
%    \operator@font for max/min) in \AtBeginDocument
%    is the only fix that held (mathastext and \DeclareSymbolFont{operators}
%    did not), and it is plain pdflatex-legal LaTeX;
%  - metric slides (2026-10-14 redesign): slide A = plans with the matching
%    parts marked (\ck, \xx, \hlt, \hlx) + ONE \bigformula + \legend + \whyline;
%    slide B = \workedmath + ONE \bigresult + ONE \blindspot callout; simple
%    metrics get one slide; edge-case conventions go in the footer only;
%  - every worked number, toy statistics included, is printed by
%    Slides/score_metric_examples.py; REAL RUN lines quote experiments/exp-005.
% If plex-sans is swapped for lmodern, delete this block.
\makeatletter
\AtBeginDocument{%
  \DeclareSymbolFont{plexdigits}{T1}{IBMPlexSans-TLF}{m}{n}%
  \SetMathAlphabet{\mathrm}{normal}{T1}{IBMPlexSans-TLF}{m}{n}%
  \DeclareMathSymbol{0}{\mathalpha}{plexdigits}{`0}\DeclareMathSymbol{1}{\mathalpha}{plexdigits}{`1}%
  \DeclareMathSymbol{2}{\mathalpha}{plexdigits}{`2}\DeclareMathSymbol{3}{\mathalpha}{plexdigits}{`3}%
  \DeclareMathSymbol{4}{\mathalpha}{plexdigits}{`4}\DeclareMathSymbol{5}{\mathalpha}{plexdigits}{`5}%
  \DeclareMathSymbol{6}{\mathalpha}{plexdigits}{`6}\DeclareMathSymbol{7}{\mathalpha}{plexdigits}{`7}%
  \DeclareMathSymbol{8}{\mathalpha}{plexdigits}{`8}\DeclareMathSymbol{9}{\mathalpha}{plexdigits}{`9}%
  \DeclareMathSymbol{.}{\mathord}{plexdigits}{`.}%
  \SetMathAlphabet{\mathbf}{normal}{T1}{IBMPlexSans-TLF}{b}{n}%
  \def\operator@font{\mathgroup\symplexdigits}}
\makeatother

% Letter-spacing for the small uppercase eyebrows.  microtype supplies
% \textls under pdflatex; the \providecommand below is a harmless no-op
% that only takes effect if it does not.  If letterspacing misbehaves,
% redefine \lsp as:  \renewcommand{\lsp}[1]{#1}
\providecommand{\textls}[2][]{#2}
\newcommand{\lsp}[1]{\textls[140]{#1}}

% --------------------------------------------------------------- colours
\definecolor{paper}{HTML}{F5F6F4}
\definecolor{sunk}{HTML}{ECEEEB}
\definecolor{ink}{HTML}{16211F}
\definecolor{inktwo}{HTML}{3C4A47}
\definecolor{muted}{HTML}{687573}
\definecolor{rulec}{HTML}{D7DDD9}
\definecolor{ruletwo}{HTML}{C3CCC7}
\definecolor{accent}{HTML}{0D6B60}
\definecolor{accentsoft}{HTML}{DDEAE7}
\definecolor{accentink}{HTML}{0A4F47}
\definecolor{gold}{HTML}{A8701A}
\definecolor{goldsoft}{HTML}{F2E7D4}
\definecolor{bad}{HTML}{A33227}
\definecolor{badsoft}{HTML}{F4DEDB}

% ----------------------------------------------------------- beamer base
\usetheme{default}
\usecolortheme{default}
\setbeamertemplate{navigation symbols}{}
\setbeamertemplate{headline}{}
\setbeamertemplate{footline}{}
\setbeamertemplate{blocks}[default]
\setbeamersize{text margin left=9mm,text margin right=9mm}
\setbeamercolor{background canvas}{bg=paper}
\setbeamercolor{normal text}{fg=ink}
\usebeamercolor[fg]{normal text}
\setbeamerfont{frametitle}{size=\normalsize}
\setlength{\parindent}{0pt}
\setlength{\parskip}{0pt}

% ------------------------------------------------------------- furniture
% eyebrow + headline block at the top of a content slide
\newcommand{\slidehead}[2]{%
  {\color{accent}\fontsize{6.8}{8}\selectfont\bfseries\lsp{#1}\par}%
  \vspace{2.4mm}%
  {\color{ink}\fontsize{15}{18}\selectfont #2\par}%
  \vspace{3.4mm}%
}
% same, but the eyebrow is grey rather than accent (HTML: .eyebrow.plain)
\newcommand{\slideheadplain}[2]{%
  {\color{muted}\fontsize{6.8}{8}\selectfont\bfseries\lsp{#1}\par}%
  \vspace{2.4mm}%
  {\color{ink}\fontsize{15}{18}\selectfont #2\par}%
  \vspace{3.4mm}%
}

\newcommand{\thinrule}[1]{\noindent\textcolor{#1}{\rule{\linewidth}{0.4pt}}\par}

% bottom source line: one item, two items, or a mono-only item
\newcommand{\srcline}[1]{%
  \vfill\thinrule{rulec}\vspace{1.5mm}%
  {\color{muted}\fontsize{6.6}{8}\selectfont #1\par}}
\newcommand{\srcpair}[2]{%
  \vfill\thinrule{rulec}\vspace{1.5mm}%
  {\color{muted}\fontsize{6.6}{8}\selectfont #1\hfill{\ttfamily #2}\par}}
\newcommand{\srcmono}[1]{%
  \vfill\thinrule{rulec}\vspace{1.5mm}%
  {\color{muted}\ttfamily\fontsize{6.6}{8}\selectfont #1\par}}

% callout: sunk panel with a coloured left rule
\newsavebox{\cobox}
\newcommand{\calloutbg}[3]{%  #1 rule colour  #2 background  #3 text
  \par\vspace{2.6mm}%
  \begin{lrbox}{\cobox}%
    \begin{minipage}{\dimexpr\linewidth-9mm\relax}%
      \color{inktwo}\footnotesize #3%
    \end{minipage}%
  \end{lrbox}%
  \noindent{\setlength{\fboxsep}{0pt}%
    \colorbox{#1}{\raisebox{0pt}[\dimexpr\ht\cobox+2.2mm\relax][\dimexpr\dp\cobox+2.2mm\relax]{\rule{1.6pt}{0pt}}}%
    \colorbox{#2}{\raisebox{0pt}[\dimexpr\ht\cobox+2.2mm\relax][\dimexpr\dp\cobox+2.2mm\relax]{\hspace{3mm}\usebox{\cobox}\hspace{3mm}}}}%
  \par\vspace{2.6mm}%
}
\newcommand{\callout}[1]{\calloutbg{accent}{sunk}{#1}}
\newcommand{\calloutwarn}[1]{\calloutbg{gold}{sunk}{#1}}
\newcommand{\calloutopen}[1]{\calloutbg{bad}{sunk}{#1}}

% a plain tinted text box (ladder rungs, quotes)
\newcommand{\txtbox}[3]{%  #1 background  #2 width  #3 content
  {\setlength{\fboxsep}{1.8mm}\colorbox{#1}{\begin{minipage}[t]{#2}#3\end{minipage}}}}

% stat: a large number above a small grey caption
\newcommand{\stat}[3]{%  #1 colour  #2 value  #3 caption
  {\color{#1}\fontsize{19}{21}\selectfont #2\par}%
  \vspace{1.6mm}%
  {\color{muted}\scriptsize #3\par}}

% "point" blocks
\newcommand{\ptag}[1]{{\color{muted}\fontsize{6.2}{7.5}\selectfont\bfseries\lsp{#1}\par}\vspace{1.4mm}}
\newcommand{\phead}[1]{{\color{ink}\fontsize{10.5}{12.5}\selectfont\bfseries #1\par}\vspace{1.3mm}}
\newcommand{\pbody}[1]{{\color{inktwo}\footnotesize #1\par}}

% label + wrapped body sharing one baseline.  Used instead of a two-column
% tabular: \arraystretch (1.25 above) shifts p-column text down relative to
% its neighbour, and \makebox + \parbox[t] cannot drift.
% Both sides are \parbox[t] and the size is set OUTSIDE both boxes: a size
% change inside a [t]-aligned box positions its first line against the outer
% font's strut, which is what makes label and body drift apart.
\newcommand{\defrow}[3][30mm]{%
  \par\noindent\footnotesize
  \parbox[t]{#1}{\raggedright\bfseries\textcolor{ink}{#2}}%
  \parbox[t]{\dimexpr\linewidth-#1\relax}{\raggedright\textcolor{inktwo}{#3}}%
  \par\vspace{2.6mm}}

% \defrow whose body may hold several paragraphs (formula lines): \color
% inside the \parbox instead of \textcolor, which cannot take a \par.
% \leavevmode first: a \color in vertical mode becomes the [t] box's first
% item and drops the body a line below its label.
\newcommand{\defrowpar}[3][30mm]{%
  \par\noindent\footnotesize
  \parbox[t]{#1}{\raggedright\bfseries\textcolor{ink}{#2}}%
  \parbox[t]{\dimexpr\linewidth-#1\relax}{\raggedright\leavevmode\color{inktwo}#3}%
  \par\vspace{2.6mm}}

% right-aligned label + sublabel beside a tinted example box.  Same rule as
% \defrow: the size is set outside both boxes so their first lines share a
% baseline.  Used by the masking and normalisation ladders.
\newcommand{\ladderrow}[5][28mm]{% width, label, sublabel, box colour, content
  \par\noindent\footnotesize
  \parbox[t]{#1}{\raggedleft\bfseries\textcolor{ink}{#2}\par
   \vspace{0.3mm}{\scriptsize\textcolor{muted}{#3}\par}}%
  \hspace{3mm}%
  \txtbox{#4}{\dimexpr\linewidth-#1-7mm\relax}{\ttfamily#5}%
  \par\vspace{1.8mm}}

% label + tinted example box, same rule: size outside, \parbox[t] both sides
\newcommand{\exrow}[4][20mm]{%
  \par\noindent\footnotesize
  \parbox[t]{#1}{\raggedright\bfseries\textcolor{ink}{#2}}%
  \txtbox{#3}{\dimexpr\linewidth-#1-4mm\relax}{\ttfamily#4}%
  \par\vspace{2.2mm}}

% table column head, gold/mono step placeholder, highlight, tick, pill
\newcommand{\colhead}[1]{{\color{accent}\fontsize{6.5}{8}\selectfont\bfseries\lsp{#1}}}
\newcommand{\plc}[1]{{\ttfamily\color{gold}#1}}
\newcommand{\hlt}[1]{{\setlength{\fboxsep}{1pt}\colorbox{accentsoft}{\textcolor{accentink}{#1}}}}
% gold plan shared by every metric-example slide (all example numbers are scored on it)
\newcommand{\xx}{\textcolor{bad}{$\times$}}
% a labelled block of stacked plan lines: label in the table-header style, then the steps
\newcommand{\planblock}[2]{\colhead{#1}\par\vspace{0.8mm}{\footnotesize\color{inktwo}#2\par}\vspace{2.4mm}}
% gold plan on the left, one predicted plan on the right, both stacked
\newcommand{\planpair}[4]{%  left label, left lines, right label, right lines
 \begin{columns}[T,onlytextwidth]
 \begin{column}{0.47\textwidth}\planblock{#1}{#2}\end{column}
 \begin{column}{0.47\textwidth}\planblock{#3}{#4}\end{column}
 \end{columns}}
\newcommand{\goldlines}{1~Who directed \textbf{The Matrix}?\par 2~Which other films did \plc{[\#1]} direct?\par 3~What languages are spoken in \plc{[\#2]}?}
% the five example plans (scored by Slides/score_metric_examples.py)
\newcommand{\revlines}{1~What languages are spoken in \plc{[\#2]}?\par 2~Which other films did \plc{[\#1]} direct?\par 3~Who directed \textbf{The Matrix}?}
\newcommand{\mergedlines}{1~Who directed \textbf{The Matrix}?\par 2~What languages are spoken in films \plc{[\#1]} directed?}
\newcommand{\unlinkedlines}{1~Who directed \textbf{The Matrix}?\par 2~Which other films did that director direct?\par 3~What languages are spoken in those films?}
\newcommand{\extralines}{1~Who directed \textbf{The Matrix}?\par 2~Which other films did \plc{[\#1]} direct?\par 3~What are the titles of those films?\par 4~What languages are spoken in \plc{[\#2]}?}
\newcommand{\punctlines}{1~Who directed \textbf{The Matrix}\hlt{.}\par 2~Which other films did \plc{[\#1]} direct?\par 3~What languages are spoken in \plc{[\#2]}?}
% HOW / WHY block: label above text, same label style as the calculation beside it
\newcommand{\hw}[2]{\colhead{#1}\par\vspace{0.8mm}{\footnotesize\color{inktwo}#2\par}\vspace{2.6mm}}
% tick mark, drawn rather than taken from a symbol font.
% If it misbehaves, replace this definition with:  \newcommand{\ck}{$\surd$}
\newcommand{\ck}{\tikz[baseline=0pt]{\draw[line width=.8pt,line cap=round,line join=round]
  (0ex,.55ex)--(.45ex,.05ex)--(1.25ex,1.15ex);}}
\newcommand{\pillon}[1]{\tikz[baseline=-.7ex]{\node[rounded corners=1.7mm,fill=accentsoft,
  text=accentink,inner xsep=2.4mm,inner ysep=1.2mm,font=\footnotesize\bfseries]{#1};}}

% ---- metric and statistics slides: idea + formula, then worked calculation
% red highlight: the part that does NOT match (\hlt marks the part that does)
\newcommand{\hlx}[1]{{\setlength{\fboxsep}{1pt}\colorbox{badsoft}{\textcolor{bad}{#1}}}}
% one plan line: optional mark (\ck, \xx or empty), step number, text
% (hanging indent, so a wrapped step stays clear of the mark and number)
\newcommand{\mline}[3]{{\hangindent=7mm\hangafter=1\noindent\makebox[4mm][l]{#1}\makebox[3mm][l]{#2}#3\par}}
\newcommand{\pline}[2]{{\hangindent=3mm\hangafter=1\noindent\makebox[3mm][l]{#1}#2\par}}
% the one large centred formula of a slide
\newcommand{\bigformula}[1]{\par\vspace{2.4mm}{\centering\Large$\displaystyle #1$\par}\vspace{2.6mm}}
% at most three short symbol items under the formula, one line
\newcommand{\legend}[1]{{\centering\footnotesize\color{muted}#1\par}\vspace{2mm}}
% WHY tag + one short line
\newcommand{\whyline}[1]{{\centering\footnotesize\textcolor{accent}{\bfseries\lsp{WHY}}\quad\textcolor{inktwo}{#1}\par}}
% one-slide metrics: a WORKED label over two short lines
\newcommand{\worked}[1]{\colhead{WORKED}\par\vspace{1mm}{\small\color{inktwo}#1\par}}
% two-slide metrics: substituted formula lines, aligned on the relation
\newcommand{\workedmath}[2][1.6]{{\renewcommand{\arraystretch}{#1}\large\color{inktwo}%
  $\begin{array}{@{}r@{\;}l@{}l@{}}#2\end{array}$\par}}
% the result as one big number over a short caption; optional size in pt
\newcommand{\bigresult}[3][30]{{\centering{\color{accent}\fontsize{#1}{#1}\selectfont #2\par}%
  \vspace{2mm}{\color{muted}\footnotesize #3\par}}}
% the slide's single callout, when it is a blind spot
\newcommand{\blindspot}[1]{\calloutwarn{\textbf{Blind spot.}\ #1}}
% a REAL RUN strip above the footer: the same test on a committed comparison
%   #1 the numbers  #2 footer.  Every REAL RUN line is the same committed comparison.
\newcommand{\srcreal}[2]{%
  \vfill{\setlength{\fboxsep}{1.4mm}\colorbox{accentsoft}{\parbox{\dimexpr\linewidth-2.8mm\relax}{%
    \footnotesize\parbox[t]{18mm}{\textcolor{accentink}{\bfseries\lsp{REAL RUN}}}%
    \parbox[t]{\dimexpr\linewidth-18mm\relax}{\leavevmode\color{inktwo}exp-005 $\cdot$ Qwen3.5-9B,
     unguided (A) vs oracle-guided (B) $\cdot$ $n=600$\par\vspace{0.8mm}\textcolor{ink}{#1}}}}}\par
  \vspace{1.6mm}\thinrule{rulec}\vspace{1.5mm}%
  {\color{muted}\ttfamily\fontsize{6.6}{8}\selectfont #2\par}}

% numbered chain bullet
\newcommand{\chainnum}[1]{\tikz[baseline=-.35ex]{\node[draw=accent,text=accent,circle,
  inner sep=0pt,minimum size=4.4mm,line width=.4pt,font=\fontsize{6.5}{7}\selectfont\bfseries]{#1};}}

% horizontal bar row, drawn inside a tikzpicture with x=1mm,y=1mm
%   #1 y  #2 label  #3 percent 0-100  #4 value  #5 fill colour  #6 track length mm
\newcommand{\barrow}[6]{%
  \node[anchor=east,font=\footnotesize,text=ink] at (13,{#1}) {#2};
  \fill[sunk] (15,{#1-1.6}) rectangle ({15+#6},{#1+1.6});
  \fill[#5] (15,{#1-1.6}) rectangle ({15+#6*#3/100},{#1+1.6});
  \node[anchor=west,font=\footnotesize,text=inktwo] at ({17+#6},{#1}) {#4};}

% right-aligned fixed-width column, so long uppercase headers wrap
% instead of forcing the table past the slide width
\newcolumntype{R}[1]{>{\raggedleft\arraybackslash}p{#1}}
\newcolumntype{S}[1]{>{\scriptsize\raggedleft\arraybackslash}p{#1}}
\renewcommand{\arraystretch}{1.25}
\setlength{\tabcolsep}{3pt}

\begin{document}

% ===================================================== 1 -- title ======
\begin{frame}[t]
\vspace{9mm}
{\color{muted}\fontsize{7.5}{9}\selectfont\bfseries\lsp{MASTER'S THESIS \textperiodcentered\ SABANCI UNIVERSITY}\par}
\vspace{6mm}
{\color{ink}\fontsize{26}{30}\selectfont
 Decomposing multi-hop questions\\into ordered atomic steps\par}
\vspace{7mm}
\begin{tikzpicture}[x=1mm,y=1mm]
  \fill[accent] (0,0) circle (1.7mm);
  \draw[accent,line width=.4pt] (0,0) circle (1.7mm);
  \draw[ruletwo,line width=.4pt] (2.4,0) -- (9.6,0);
  \draw[accent,line width=.4pt] (12,0) circle (1.7mm);
  \draw[ruletwo,line width=.4pt] (14.4,0) -- (21.6,0);
  \draw[accent,line width=.4pt] (24,0) circle (1.7mm);
\end{tikzpicture}
\vspace{7mm}

{\color{inktwo}\fontsize{9.5}{15}\selectfont
 {\color{ink}\bfseries Jahid Ahmadi}\\
 Supervisor: Ahmmad O. M. Saleh\\
 {\color{muted}Progress review $\cdot$ September 2026}\par}
\end{frame}

% ================================================ 2 -- motivation ======
\begin{frame}[t]
\slidehead{MOTIVATION}{Most questions worth asking are not one lookup}
\begin{columns}[T,onlytextwidth]
\begin{column}{0.53\textwidth}
 {\color{inktwo}\small
  A knowledge base holds facts as single edges. A question like the one on the right
  touches three of them, and no single lookup answers it.\par}
 \vspace{3mm}
 {\color{inktwo}\small
  Systems that answer it well do so by \hlt{planning first} --- turning one hard
  question into a short sequence of easy ones.\par}
\end{column}
\begin{column}{0.43\textwidth}
 \vspace{-2.6mm}
 \calloutbg{accent}{sunk}{\normalsize ``What languages are spoken in films that share a
  director with \textbf{The Matrix}?''}
\end{column}
\end{columns}
\end{frame}

% ================================================= 3 -- the chain ======
\begin{frame}[t]
\slidehead{MOTIVATION}{The plan is three ordered steps, and each needs the last}
\begin{tabular}{@{}p{6mm}@{\hspace{2mm}}p{118mm}@{}}
\chainnum{1} & \textcolor{ink}{\fontsize{10.5}{13}\selectfont Who directed \textbf{The Matrix}?}\\[2.2mm]
\chainnum{2} & \textcolor{ink}{\fontsize{10.5}{13}\selectfont Which other films did \plc{[\#1]} direct?}\\[2.2mm]
\chainnum{3} & \textcolor{ink}{\fontsize{10.5}{13}\selectfont What languages are spoken in \plc{[\#2]}?}
\end{tabular}
\srcline{Ordered $\cdot$ atomic $\cdot$ explicitly linked}
\end{frame}

% ======================================= 4 -- problem definition ======
\begin{frame}[t]
\slidehead{PROBLEM DEFINITION}{What goes in, what must come out}
\vspace{4mm}
\footnotesize
\begin{tabular}{@{}l@{\hspace{16mm}}l@{}}
\multicolumn{2}{@{}l}{\textbf{Input}}\\
\midrule
{\ttfamily Q}\quad a multi-hop question & \textcolor{inktwo}{natural language}\\
{\ttfamily D}\quad a domain           & \textcolor{inktwo}{the knowledge source}\\[2.5mm]
\multicolumn{2}{@{}l}{\textbf{Output}}\\
\midrule
{\ttfamily Q' =}\,$\langle q_1 \dots q_n \rangle$ & \textcolor{inktwo}{ordered sub-questions}
\end{tabular}
\par\vspace{5mm}
{\color{muted}\scriptsize
 Each $q_i$ is answerable by a single lookup in {\ttfamily D}, and may reference
 \plc{[\#j]} for $j < i$.\par}
\srcline{n is not given at inference time --- the model must decide the length too.}
\end{frame}

% =============================================== 5 -- challenges ======
\begin{frame}[t]
\slidehead{CHALLENGES}{Three ways a decomposition goes wrong}
{\color{inktwo}\small
 The three failures are independent, and one accuracy number hides all of them.\par}
\vspace{7mm}
\defrow[34mm]{Wrong length}{The model has to infer how many steps the question needs.
 Too few drop a hop; too many invent one.}
\defrow[34mm]{Broken chain}{A later step has to consume an earlier answer. Steps that
 restate the original question look fine and compose to nothing.}
\defrow[34mm]{Invented entities}{The model names people and films that are not in the
 knowledge source, so the plan cannot execute.}
\vspace{4mm}
\calloutwarn{This is why the evaluation separates \textbf{length}, \textbf{order} and
 \textbf{content} rather than reporting a single score.}
\srcline{Each failure mode maps to a metric family later in this deck.}
\end{frame}

% ==================================== 6 -- methodology overview ======
\begin{frame}[t]
\slidehead{METHODOLOGY \textperiodcentered\ OVERVIEW}{Everything that moves sits in how the prompt is built}
\vspace{3mm}
\begin{center}
\begin{tikzpicture}[x=1mm,y=1mm]
  % stage A
  \fill[sunk] (0,0) rectangle (66,26);
  \draw[accent,line width=.6pt] (0,0) rectangle (66,26);
  \node[anchor=west,text=accent,font=\ttfamily\fontsize{5}{6}\selectfont] at (0,28.6)
    {STAGE A $\cdot$ THE STUDY};
  \node[anchor=west,text=ink,font=\fontsize{7.5}{9}\selectfont\bfseries] at (3,20.5)
    {Prompt construction};
  \node[anchor=west,text=muted,font=\fontsize{5.4}{6.4}\selectfont] at (3,15.5)
    {Typed masking $\to$ few-shot pool $\to$ retrieval $\to$ rerank};
  \node[anchor=west,rounded corners=1.4mm,fill=accentsoft,text=accentink,inner xsep=1.8mm,
    inner ysep=1.1mm,font=\ttfamily\fontsize{5}{6}\selectfont] at (3,7) {masking};
  \node[anchor=west,rounded corners=1.4mm,fill=accentsoft,text=accentink,inner xsep=1.8mm,
    inner ysep=1.1mm,font=\ttfamily\fontsize{5}{6}\selectfont] at (18,7) {pool};
  \node[anchor=west,rounded corners=1.4mm,fill=accentsoft,text=accentink,inner xsep=1.8mm,
    inner ysep=1.1mm,font=\ttfamily\fontsize{5}{6}\selectfont] at (29,7) {retrieval};
  \node[anchor=west,rounded corners=1.4mm,fill=accentsoft,text=accentink,inner xsep=1.8mm,
    inner ysep=1.1mm,font=\ttfamily\fontsize{5}{6}\selectfont] at (44,7) {rerank};
  \draw[ruletwo,line width=.6pt,->] (67.5,13) -- (72.5,13);
  % stage B
  \fill[sunk] (74,0) rectangle (101,26);
  \draw[ruletwo,line width=.5pt] (74,0) rectangle (101,26);
  \node[anchor=west,text=accent,font=\ttfamily\fontsize{5}{6}\selectfont] at (74,28.6)
    {STAGE B $\cdot$ FIXED};
  \node[anchor=west,text=ink,font=\fontsize{7.5}{9}\selectfont\bfseries] at (76.5,20.5)
    {Decomposer};
  \node[anchor=west,text=muted,font=\fontsize{5.4}{6.4}\selectfont] at (76.5,15.5)
    {One open model, 7--14B};
  \node[anchor=west,text=muted,font=\fontsize{5.4}{6.4}\selectfont] at (76.5,11)
    {One decoding config};
  \node[anchor=west,text=muted,font=\fontsize{5.4}{6.4}\selectfont] at (76.5,6.5)
    {One evaluation set};
  \draw[ruletwo,line width=.6pt,->] (102.5,13) -- (107.5,13);
  % evaluation
  \fill[sunk] (109,0) rectangle (140,26);
  \draw[ruletwo,line width=.5pt] (109,0) rectangle (140,26);
  \node[anchor=west,text=accent,font=\ttfamily\fontsize{5}{6}\selectfont] at (109,28.6)
    {EVALUATION};
  \node[anchor=west,text=ink,font=\fontsize{7.5}{9}\selectfont\bfseries] at (111.5,20.5)
    {Score the plan};
  \node[anchor=west,text=muted,font=\fontsize{5.4}{6.4}\selectfont] at (111.5,15.5)
    {Against gold steps, or};
  \node[anchor=west,text=muted,font=\fontsize{5.4}{6.4}\selectfont] at (111.5,11)
    {compile and run it on};
  \node[anchor=west,text=muted,font=\fontsize{5.4}{6.4}\selectfont] at (111.5,6.5)
    {the graph};
\end{tikzpicture}
\end{center}
\callout{Stage B is held constant on purpose. If the decomposer and its decoding never change,
 then any difference between runs is attributable to \textbf{prompt construction alone}.}
\srcline{Each stage-A component gets its own slide next.}
\end{frame}

% ==================================== 7a -- masking: definition ======
\begin{frame}[t]
\slidehead{METHODOLOGY \textperiodcentered\ COMPONENT 1}{What masking is, and the two places it is applied}
{\color{inktwo}\small
 Masking replaces a named entity with a token for its \emph{type}, so questions are compared
 on shape, not on names.\par}
\vspace{3mm}
\exrow{raw}{sunk}{\textcolor{inktwo}{who directed Baby Face}}
\exrow{masked}{accentsoft}{\textcolor{accentink}{who directed }\plc{[MOVIE]}}
\vspace{3mm}
\defrow[32mm]{On the pool}{Masked once, in advance. The masked pool is a maintained artifact,
 never re-masked while a query is served.}
\defrow[32mm]{On the question}{Masked at lookup time, on every question, before it meets the pool.}
\srcmono{ADR 0003 $\cdot$ src/entity\_masking.py $\cdot$ placeholders: [MOVIE] and [PERSON]}
\end{frame}

% ====================================== 7b -- masking: variants ======
\begin{frame}[t]
\slidehead{METHODOLOGY \textperiodcentered\ COMPONENT 1}{Masking makes similarity a statement about \emph{shape}}
{\color{inktwo}\small
 Retrieval over raw text finds examples that share names. Shared names say nothing about
 how many hops a question needs.\par}
\vspace{4mm}
\ladderrow[24mm]{raw}{entities intact}{sunk}%
 {\textcolor{inktwo}{which films did Christina Ricci star in after The Addams Family in 1991}}
\ladderrow[24mm]{uniform}{one shared token}{sunk}%
 {\textcolor{inktwo}{which films did \textcolor{accent}{[MASK]} star in after \textcolor{accent}{[MASK]} in 1991}}
\ladderrow[24mm]{typed}{type label kept}{accentsoft}%
 {\textcolor{accentink}{which films did \plc{[PERSON]} star in after \plc{[MOVIE]} in 1991}}
\srcline{The pool is masked once in advance, never at query time --- dynamic masking corrupts
 titles that collide with common words.}
\end{frame}

% ===================================================== 8 -- pool ======
\begin{frame}[t]
\slidehead{METHODOLOGY \textperiodcentered\ COMPONENT 2}{The few-shot pool is a sample, and how it is drawn is a choice}
{\color{inktwo}\small
 Every prompt draws its examples from one pool. The pool is built once, before any question
 is asked --- and three strategies build it differently.\par}
\vspace{5.5mm}
\defrow[30mm]{Imbalanced}{Draws uniformly. The pool inherits whatever mix the corpus already has.}
\defrow[30mm]{Balanced}{Fills an equal quota from each hop bucket, so no hop count dominates.}
\defrow[30mm]{Clustered}{k-means over the pool, one example per centroid --- covers the spread
 rather than the frequency.}
\vspace{4mm}
\srcmono{configs/pool\_sweep.json $\cdot$ MuSiQue-Ans v1.0 train}
\end{frame}

% ======================================== 9 -- retrieval funnel ======
\begin{frame}[t]
\slidehead{METHODOLOGY \textperiodcentered\ COMPONENT 3}{Retrieval is a funnel, and the last stage is expensive}
\vspace{2mm}
\begin{center}
\begin{tikzpicture}[x=1mm,y=1mm]
  % pool
  \fill[sunk] (0,0) rectangle (32,26);
  \draw[ruletwo,line width=.5pt] (0,0) rectangle (32,26);
  \node[anchor=west,text=accent,font=\ttfamily\fontsize{5}{6}\selectfont] at (0,28.6) {POOL};
  \node[anchor=west,text=ink,font=\fontsize{7.5}{9}\selectfont\bfseries] at (2.5,21) {1,000 -- 8,000};
  \node[anchor=west,text=muted,font=\fontsize{5.4}{6.4}\selectfont] at (2.5,16.6) {masked examples,};
  \node[anchor=west,text=muted,font=\fontsize{5.4}{6.4}\selectfont] at (2.5,12.8) {embedded once};
  \foreach \r in {3.2,5.6,8} {\foreach \i in {0,...,17}
    {\fill[ruletwo] ({2.5+\i*1.6},{\r}) rectangle +(0.9,0.9);}}
  \draw[ruletwo,line width=.6pt,->] (33.5,13) -- (40.5,13);
  \node[anchor=south west,text=muted,font=\fontsize{5.4}{6.4}\selectfont] at (33.5,14) {bi-encoder};
  % top 20
  \fill[sunk] (42,2) rectangle (70,24);
  \draw[ruletwo,line width=.5pt] (42,2) rectangle (70,24);
  \node[anchor=west,text=accent,font=\ttfamily\fontsize{5}{6}\selectfont] at (42,26.6) {TOP 20};
  \node[anchor=west,text=ink,font=\fontsize{7.5}{9}\selectfont\bfseries] at (44.5,19.5) {Vector search};
  \node[anchor=west,text=muted,font=\fontsize{5.4}{6.4}\selectfont] at (44.5,15.2) {Question and pool item};
  \node[anchor=west,text=muted,font=\fontsize{5.4}{6.4}\selectfont] at (44.5,11.9) {embedded separately ---};
  \node[anchor=west,text=muted,font=\fontsize{5.4}{6.4}\selectfont] at (44.5,8.6) {fast, but they never};
  \node[anchor=west,text=muted,font=\fontsize{5.4}{6.4}\selectfont] at (44.5,5.3) {see each other.};
  \draw[ruletwo,line width=.6pt,->] (71.5,13) -- (78.5,13);
  \node[anchor=south west,text=muted,font=\fontsize{5.4}{6.4}\selectfont] at (71.5,14) {rerank};
  % top 5
  \fill[accentsoft] (80,4) rectangle (108,22);
  \draw[accent,line width=.6pt] (80,4) rectangle (108,22);
  \node[anchor=west,text=accent,font=\ttfamily\fontsize{5}{6}\selectfont] at (80,24.6) {TOP 5};
  \node[anchor=west,text=ink,font=\fontsize{7.5}{9}\selectfont\bfseries] at (82.5,18) {Cross-encoder};
  \node[anchor=west,text=muted,font=\fontsize{5.4}{6.4}\selectfont] at (82.5,13.7) {Scored jointly with};
  \node[anchor=west,text=muted,font=\fontsize{5.4}{6.4}\selectfont] at (82.5,10.4) {the question. Accurate,};
  \node[anchor=west,text=muted,font=\fontsize{5.4}{6.4}\selectfont] at (82.5,7.1) {far slower.};
  \draw[ruletwo,line width=.6pt,->] (109.5,13) -- (116.5,13);
  % prompt
  \fill[sunk] (118,4) rectangle (140,22);
  \draw[ruletwo,line width=.5pt] (118,4) rectangle (140,22);
  \node[anchor=west,text=accent,font=\ttfamily\fontsize{5}{6}\selectfont] at (118,24.6) {PROMPT};
  \node[anchor=west,text=ink,font=\fontsize{7.5}{9}\selectfont\bfseries] at (120.5,18) {5 examples};
  \node[anchor=west,text=muted,font=\fontsize{5.4}{6.4}\selectfont] at (120.5,13.7) {placed highest-};
  \node[anchor=west,text=muted,font=\fontsize{5.4}{6.4}\selectfont] at (120.5,10.4) {similarity first};
\end{tikzpicture}
\end{center}
\srcmono{bi-encoder e5-small-v2 $\cdot$ cross-encoder rerank over top-20}
\end{frame}

% ==================================== 11 -- research questions ======
\begin{frame}[t]
\slidehead{RESEARCH QUESTIONS}{Now that the method is on the table, these are the questions it asks}
\vspace{1mm}
\begin{tabular}{@{}p{6mm}@{\hspace{2mm}}p{126mm}@{}}
\chainnum{1} & \textcolor{ink}{\fontsize{10}{12.5}\selectfont
 Does \textbf{how the few-shot examples are chosen} --- masking, pool, retrieval, rerank ---
 determine whether a mid-sized model can decompose at all?}\\[3.6mm]
\chainnum{2} & \textcolor{ink}{\fontsize{10}{12.5}\selectfont
 Does \textbf{knowing the hop count in advance} help, and is the answer the same across models?}\\[3.6mm]
\chainnum{3} & \textcolor{ink}{\fontsize{10}{12.5}\selectfont
 Does \textbf{fine-tuning} beat prompting --- and does it still hold at a depth it never saw
 in training?}\\[3.6mm]
\chainnum{4} & \textcolor{ink}{\fontsize{10}{12.5}\selectfont
 Do \textbf{better plans produce better answers} end to end?}
\end{tabular}
\srcline{Every claim that follows is paired to one of these four.}
\end{frame}

% ========================================= 12 -- divider: setup ======
{\setbeamercolor{background canvas}{bg=sunk}
\begin{frame}[t]
\vspace{18mm}
{\color{accent}\fontsize{11}{13}\selectfont Part two\par}
\vspace{4mm}
{\color{ink}\fontsize{24}{28}\selectfont Experimental setup\par}
\vspace{5mm}
{\color{inktwo}\fontsize{11}{15}\selectfont
 Datasets, what counts as a good decomposition, and how a difference is\\declared real.\par}
\end{frame}
}

% ================================================= 13 -- datasets ======
\begin{frame}[t]
\slidehead{SETUP \textperiodcentered\ DATA}{Two datasets doing two different jobs}
\begin{columns}[T,onlytextwidth]
\begin{column}{0.47\textwidth}
\footnotesize
\begin{tabular}{@{}p{14mm}R{16mm}R{18mm}R{16mm}@{}}
\colhead{MUSIQUE} & \colhead{TRAIN} & \colhead{DEV AVAIL.} & \colhead{DEV USED}\\
\midrule
2-hop & \textcolor{inktwo}{14,376} & \textcolor{inktwo}{1,252} & \textcolor{inktwo}{200}\\
3-hop & \textcolor{inktwo}{4,387} & \textcolor{inktwo}{760} & \textcolor{inktwo}{200}\\
4-hop & \textcolor{inktwo}{1,175} & \textcolor{inktwo}{405} & \textcolor{inktwo}{200}\\
\bottomrule
\end{tabular}
\par\vspace{3.5mm}
{\color{inktwo}\scriptsize
 Ships \textbf{human-written gold decompositions} --- so the plan itself can be scored.
 Carries all decomposition-quality evaluation.\par}
\end{column}
\begin{column}{0.47\textwidth}
\footnotesize
\begin{tabular}{@{}p{16mm}R{20mm}R{18mm}@{}}
\colhead{METAQA} & \colhead{KB QUESTIONS} & \colhead{EVAL USED}\\
\midrule
1-hop & \textcolor{inktwo}{96,106} & \textcolor{inktwo}{100}\\
2-hop & \textcolor{inktwo}{118,980} & \textcolor{inktwo}{100}\\
3-hop & \textcolor{inktwo}{114,196} & \textcolor{inktwo}{100}\\
\bottomrule
\end{tabular}
\par\vspace{3.5mm}
{\color{inktwo}\scriptsize
 A movie-domain triple store with \textbf{no gold decompositions} --- so a plan is judged
 only by whether it executes and returns the right answer.\par}
\end{column}
\end{columns}
\end{frame}

% ================ 14-16 -- metrics block: overview x2, normalisation, example, metric slides ======
\begin{frame}[t]
\slidehead{SETUP \textperiodcentered\ METRICS}{The report card: six numbers, never blended}
\footnotesize
\begin{tabular}{@{}p{34mm}p{86mm}R{14mm}@{}}
\colhead{METRIC} & \colhead{WHAT IT CHECKS} & \colhead{BETTER}\\
\midrule
\textbf{Break EM} & \textcolor{inktwo}{The whole plan is identical to gold.} & $\uparrow$\\
\textbf{SARI} & \textcolor{inktwo}{Edits the question the way gold does.} & $\uparrow$\\
\textbf{GED} & \textcolor{inktwo}{Graph distance to gold: wording, links, length.} & $\downarrow$\\
\textbf{Chain validity} & \textcolor{inktwo}{Every \plc{[\#k]} points to an earlier step.} & $\uparrow$\\
\textbf{Hop-count EM} & \textcolor{inktwo}{As many steps as gold.} & $\uparrow$\\
\textbf{Under / over rate} & \textcolor{inktwo}{Too few or too many steps, as a pair.} & $\downarrow$\\
\bottomrule
\end{tabular}
\srcmono{ADR 0029 $\cdot$ a plan can pass some and fail others, so all six are reported}
\end{frame}

\begin{frame}[t]
\slidehead{SETUP \textperiodcentered\ METRICS}{Also reported, and the end-to-end pair}
\footnotesize
\begin{tabular}{@{}p{34mm}p{86mm}R{14mm}@{}}
\colhead{PLAN QUALITY} & \colhead{WHAT IT CHECKS} & \colhead{BETTER}\\
\midrule
\textbf{Exact match} & \textcolor{inktwo}{Every step identical to gold.} & $\uparrow$\\
\textbf{Step F1} & \textcolor{inktwo}{The right steps are present; order ignored.} & $\uparrow$\\
\textbf{Ordered step acc.} & \textcolor{inktwo}{The right step in the right place.} & $\uparrow$\\
\textbf{ROUGE-L F1} & \textcolor{inktwo}{Longest common word sequence with gold.} & $\uparrow$\\
\midrule
\colhead{END TO END} & & \\
\midrule
\textbf{Answer EM} & \textcolor{inktwo}{The final answer equals gold's.} & $\uparrow$\\
\textbf{Answer F1} & \textcolor{inktwo}{Word overlap with gold's final answer.} & $\uparrow$\\
\bottomrule
\end{tabular}
\srcline{The old weighted composite is legacy: it ranks nothing here.}
\end{frame}

\begin{frame}[t]
\slidehead{SETUP \textperiodcentered\ METRICS}{Normalisation $\nu$: what it does, where it is used}
\ladderrow{as generated}{raw step}{sunk}%
 {\textcolor{inktwo}{Which well-known films did \plc{[\#2]} direct?}}
\ladderrow{1 $\cdot$ lower}{+ strip outer space}{sunk}%
 {\textcolor{inktwo}{which well-known films did \plc{[\#2]} direct?}}
\ladderrow{2 $\cdot$ delete}{all but word chars, space, \#}{sunk}%
 {\textcolor{inktwo}{which wellknown films did \plc{\#2} direct}}
\ladderrow{3 $\cdot$ collapse}{whitespace runs}{accentsoft}%
 {\textcolor{accentink}{which wellknown films did \plc{\#2} direct}}
\vspace{2.5mm}
\begin{columns}[T,onlytextwidth]
\begin{column}{0.47\textwidth}
 \ptag{USED BY}
 \pbody{Exact match, step F1, ordered step accuracy. The rest have their own rules.}
\end{column}
\begin{column}{0.47\textwidth}
 \ptag{WHY}
 \pbody{A capital letter, or \plc{[\#2]} vs \plc{\#2}, is never a wrong step.}
\end{column}
\end{columns}
\end{frame}

\begin{frame}[t]
\slidehead{SETUP \textperiodcentered\ METRICS}{The example every metric is scored on}
\begin{columns}[T,onlytextwidth]
\begin{column}{0.47\textwidth}
\planblock{GOLD PLAN}{\goldlines}
\planblock{MERGED: STEPS 2 AND 3 AS ONE}{\mergedlines}
\planblock{PUNCTUATION ONLY: STEP 1 ENDS IN A FULL STOP}{\punctlines}
\end{column}
\begin{column}{0.47\textwidth}
\planblock{REVERSED: STEPS 3, 2, 1}{\revlines}
\planblock{EXTRA STEP ADDED}{\extralines}
\planblock{UNLINKED: NO \plc{[\#k]} LINKS}{\unlinkedlines}
\end{column}
\end{columns}
\srcline{Each metric slide scores one of these plans against the gold plan.}
\end{frame}

% ---- per-metric slides.  Every number below is printed by
% ---- Slides/score_metric_examples.py, which calls the evaluator's own score_item.
% ---- Pattern: slide A = plans with the matching parts marked + ONE formula + legend + WHY;
% ---- slide B = the formula with numbers substituted + ONE big result + ONE blind-spot callout.
% ---- Edge-case conventions live in the footer, never in the body.

% ------------------------------------------------------------ exact match (one slide)
\begin{frame}[t]
\slidehead{SETUP \textperiodcentered\ METRICS}{Exact match: is every step identical?}
\footnotesize
\begin{tabular}{@{}c@{\hspace{3mm}}l@{\hspace{8mm}}l@{\hspace{8mm}}c@{}}
\colhead{$i$} & \colhead{GOLD} & \colhead{REVERSED} & \\
\midrule
1 & \textcolor{inktwo}{Who directed \textbf{The Matrix}?} & \textcolor{inktwo}{What languages are spoken in \plc{[\#2]}?} & \xx\\
2 & \textcolor{inktwo}{Which other films did \plc{[\#1]} direct?} & \textcolor{inktwo}{Which other films did \plc{[\#1]} direct?} & \ck\\
3 & \textcolor{inktwo}{What languages are spoken in \plc{[\#2]}?} & \textcolor{inktwo}{Who directed \textbf{The Matrix}?} & \xx\\
\bottomrule
\end{tabular}
\bigformula{\mathrm{EM}=\mathbf{1}\big[\,|p|=|g|\ \wedge\ \nu(p_i)=\nu(g_i)\ \ \forall i\,\big]}
\begin{columns}[c,onlytextwidth]
\begin{column}{0.62\textwidth}
\worked{$|p|=|g|=3$\ \ck,\quad but $i=1$ and $i=3$\ \xx\par\vspace{1mm}
 Punctuation-only plan: $\mathrm{EM}=1$, since $\nu$ strips the mark}
\end{column}
\begin{column}{0.34\textwidth}
\bigresult{0}{EM, reversed plan}
\end{column}
\end{columns}
\srcmono{exact\_match\_rate $\cdot$ $p,g$: predicted, gold steps $\cdot$ run score: mean over $N$ items}
\end{frame}

% ------------------------------------------------------------ step F1 (two slides)
\begin{frame}[t]
\slidehead{SETUP \textperiodcentered\ METRICS}{Step F1: are the right steps there?}
\planpair{GOLD PLAN}{\mline{\ck}{1}{Who directed \textbf{The Matrix}?}%
 \mline{\xx}{2}{Which other films did \plc{[\#1]} direct?}%
 \mline{\xx}{3}{What languages are spoken in \plc{[\#2]}?}}%
 {MERGED PLAN}{\mline{\ck}{1}{Who directed \textbf{The Matrix}?}%
 \mline{\xx}{2}{What languages are spoken in films \plc{[\#1]} directed?}}
\bigformula{F_1=\frac{2\,|S_p\cap S_g|}{|S_p|+|S_g|}}
\legend{$S_p,S_g$: sets of normalised steps\qquad \ck\ in both sets\qquad same as $2PR/(P+R)$}
\whyline{Partial credit for a plan that is mostly right.}
\srcmono{step\_f1\_macro $\cdot$ run score: mean over $N$ items}
\end{frame}

\begin{frame}[t]
\slidehead{SETUP \textperiodcentered\ METRICS}{Step F1, worked on the merged plan}
\vspace{3mm}
\begin{columns}[c,onlytextwidth]
\begin{column}{0.60\textwidth}
\colhead{WORKED $\cdot$ MERGED}\par\vspace{1mm}
\workedmath{
 |S_p\cap S_g| &{}= 1 &\quad \text{(step 1)}\\
 |S_p|,\ |S_g| &{}= 2,\ 3\\
 P,\ R &{}= \tfrac12,\ \tfrac13\\
 F_1 &{}= \dfrac{2\cdot1}{2+3}}
\end{column}
\begin{column}{0.36\textwidth}
\bigresult{0.40}{step F1, merged plan}
\end{column}
\end{columns}
\blindspot{Reversed scores 1.00. A set cannot see order.}
\srcmono{step\_f1 $\cdot$ $F_1=0$ when nothing overlaps; 1 when both plans are empty}
\end{frame}

% ------------------------------------------------------------ ordered step accuracy (two slides)
\begin{frame}[t]
\slidehead{SETUP \textperiodcentered\ METRICS}{Ordered step accuracy: right step, right place?}
\footnotesize
\begin{tabular}{@{}c@{\hspace{3mm}}l@{\hspace{8mm}}l@{\hspace{8mm}}c@{}}
\colhead{$i$} & \colhead{GOLD} & \colhead{EXTRA STEP ADDED} & \\
\midrule
1 & \textcolor{inktwo}{Who directed \textbf{The Matrix}?} & \textcolor{inktwo}{Who directed \textbf{The Matrix}?} & \ck\\
2 & \textcolor{inktwo}{Which other films did \plc{[\#1]} direct?} & \textcolor{inktwo}{Which other films did \plc{[\#1]} direct?} & \ck\\
3 & \textcolor{inktwo}{What languages are spoken in \plc{[\#2]}?} & \textcolor{inktwo}{What are the titles of those films?} & \xx\\
4 & \textcolor{muted}{---} & \textcolor{inktwo}{What languages are spoken in \plc{[\#2]}?} & \textcolor{muted}{--}\\
\bottomrule
\end{tabular}
\bigformula{\mathrm{OA}=\frac{1}{\max(|p|,|g|)}\sum_{i=1}^{\min(|p|,|g|)}\mathbf{1}\big[\nu(p_i)=\nu(g_i)\big]}
\legend{position $i$ against position $i$\qquad $|x|$: number of steps}
\whyline{Catches wrong order and wrong length.}
\srcmono{ordered\_step\_accuracy\_macro $\cdot$ run score: mean over $N$ items}
\end{frame}

\begin{frame}[t]
\slidehead{SETUP \textperiodcentered\ METRICS}{Ordered step accuracy, worked on the extra-step plan}
\vspace{3mm}
\begin{columns}[c,onlytextwidth]
\begin{column}{0.60\textwidth}
\colhead{WORKED $\cdot$ EXTRA STEP}\par\vspace{1mm}
\workedmath{
 \text{matches} &{}= 2 &\quad (i=1,2)\\
 \max(|p|,|g|) &{}= \max(4,3)=4\\
 \mathrm{OA} &{}= 2\,/\,4}
\end{column}
\begin{column}{0.36\textwidth}
\bigresult{0.50}{ordered step accuracy, extra step}
\end{column}
\end{columns}
\blindspot{Reversed holds all three gold steps, yet scores 0.33.}
\srcmono{ordered\_step\_accuracy $\cdot$ merged and unlinked also score 0.33}
\end{frame}

% ------------------------------------------------------------ ROUGE-L (two slides)
\begin{frame}[t]
\slidehead{SETUP \textperiodcentered\ METRICS}{ROUGE-L F1: how much wording overlaps?}
\planpair{GOLD PLAN}{\pline{1}{Who directed \textbf{The Matrix}?}%
 \pline{2}{Which other films did \hlx{[\#1]} direct?}%
 \pline{3}{What languages are spoken in \hlx{[\#2]?}}}%
 {UNLINKED PLAN}{\pline{1}{Who directed \textbf{The Matrix}?}%
 \pline{2}{Which other films did \hlx{that director} direct?}%
 \pline{3}{What languages are spoken in \hlx{those films?}}}
\bigformula{F=\frac{2\,|\mathrm{LCS}|}{|t(p)|+|t(g)|}}
\legend{$t(x)$: all of a plan's words\qquad LCS: longest common word subsequence\qquad \hlx{red}: outside it}
\whyline{Credit for near-miss wording inside a step.}
\srcmono{rouge\_l\_f1\_macro $\cdot$ own tokeniser, not $\nu$: lowercased, split on spaces, punctuation kept}
\end{frame}

\begin{frame}[t]
\slidehead{SETUP \textperiodcentered\ METRICS}{ROUGE-L F1, worked on the unlinked plan}
\vspace{3mm}
\begin{columns}[c,onlytextwidth]
\begin{column}{0.60\textwidth}
\colhead{WORKED $\cdot$ UNLINKED}\par\vspace{1mm}
\workedmath{
 |t(p)|,\ |t(g)| &{}= 18,\ 16\\
 |\mathrm{LCS}| &{}= 14\\
 P,\ R &{}= \tfrac{14}{18},\ \tfrac{14}{16} = 0.778,\ 0.875\\
 F &{}= \dfrac{2\cdot14}{18+16}}
\vspace{2.5mm}{\footnotesize\color{muted}Plan words outside the LCS: \hlx{that director}, \hlx{those films?}\par}
\end{column}
\begin{column}{0.36\textwidth}
\bigresult{0.82}{ROUGE-L F1, unlinked plan}
\end{column}
\end{columns}
\blindspot{Reversed has every gold word, yet scores 0.38.}
\srcmono{rouge\_l\_f1 $\cdot$ $F=2PR/(P+R)$ with $P=|\mathrm{LCS}|/|t(p)|$, $R=|\mathrm{LCS}|/|t(g)|$}
\end{frame}

% ------------------------------------------------------------ hop count (one slide)
\begin{frame}[t]
\slidehead{SETUP \textperiodcentered\ METRICS}{Hop-count EM and under/over: is the length right?}
\planpair{GOLD PLAN $\cdot$ 3 STEPS}{\goldlines}{MERGED PLAN $\cdot$ 2 STEPS}{\pline{1}{Who directed \textbf{The Matrix}?}%
 \pline{2}{What languages are spoken in films \plc{[\#1]} directed?}}
\bigformula{\mathrm{HEM}=\mathbf{1}\big[|p|=|g|\big]\qquad \mathrm{U}=\mathbf{1}\big[|p|<|g|\big]\qquad \mathrm{O}=\mathbf{1}\big[|p|>|g|\big]}
\begin{columns}[c,onlytextwidth]
\begin{column}{0.62\textwidth}
\worked{$|p|=2<|g|=3$\par\vspace{1mm}
 Reversed: $\mathrm{HEM}=1$. Right length, broken plan.}
\end{column}
\begin{column}{0.34\textwidth}
\bigresult{0\,/\,1\,/\,0}{HEM / U / O, merged plan}
\end{column}
\end{columns}
\srcmono{hop\_count\_exact\_match\_rate $\cdot$ under / over rates never summed: too few is worse $\cdot$ ADR 0017}
\end{frame}

% ------------------------------------------------------------ chain validity (two slides)
\begin{frame}[t]
\slidehead{SETUP \textperiodcentered\ METRICS}{Chain validity: do the steps link up?}
\planpair{GOLD PLAN}{\mline{}{1}{Who directed \textbf{The Matrix}?}%
 \mline{\ck}{2}{Which other films did \hlt{[\#1]} direct?}%
 \mline{\ck}{3}{What languages are spoken in \hlt{[\#2]}?}}%
 {REVERSED PLAN}{\mline{\xx}{1}{What languages are spoken in \hlx{[\#2]}?}%
 \mline{\ck}{2}{Which other films did \hlt{[\#1]} direct?}%
 \mline{}{3}{Who directed \textbf{The Matrix}?}}
\bigformula{\mathrm{CV}=\begin{cases}1 & L_g=0\\ 0 & L_g>0,\ L_p=0\\ V_p\,/\,L_p & \text{otherwise}\end{cases}}
\legend{$L_x$: \plc{\#k} links in plan $x$\qquad $V_p$: links in step $i$ with $1\le k<i$}
\whyline{Links make it a plan. No free credit for silence.}
\srcmono{chain\_validity\_macro $\cdot$ run score: mean over $N$ items $\cdot$ ADR 0026}
\end{frame}

\begin{frame}[t]
\slidehead{SETUP \textperiodcentered\ METRICS}{Chain validity, worked on the reversed plan}
\vspace{3mm}
\begin{columns}[c,onlytextwidth]
\begin{column}{0.60\textwidth}
\colhead{WORKED $\cdot$ REVERSED}\par\vspace{1mm}
\workedmath{
 \text{step 1} \to \#2\text{:} &\ 2\not<1 &\quad \text{\xx}\\
 \text{step 2} \to \#1\text{:} &\ 1<2 &\quad \text{\ck}\\
 \mathrm{CV} &{}= V_p\,/\,L_p = 1\,/\,2\\
 \text{unlinked plan} &\ L_p=0\ \Rightarrow\ \mathrm{CV}=0}
\end{column}
\begin{column}{0.36\textwidth}
\bigresult{0.50}{chain validity, reversed plan}
\end{column}
\end{columns}
\blindspot{Merged scores 1.00: links must point back, not match gold.}
\srcmono{chain\_validity $\cdot$ links counted by \plc{\#k}, with or without brackets}
\end{frame}

% ------------------------------------------------------------ Break EM (one slide)
\begin{frame}[t]
\slidehead{SETUP \textperiodcentered\ METRICS}{Break EM: is the whole plan string identical?}
\planpair{GOLD PLAN}{\pline{1}{Who directed \textbf{The Matrix}\hlt{?}}%
 \pline{2}{Which other films did \plc{[\#1]} direct?}%
 \pline{3}{What languages are spoken in \plc{[\#2]}?}}%
 {PUNCTUATION-ONLY PLAN}{\pline{1}{Who directed \textbf{The Matrix}\hlx{.}}%
 \pline{2}{Which other films did \plc{[\#1]} direct?}%
 \pline{3}{What languages are spoken in \plc{[\#2]}?}}
\bigformula{\mathrm{BEM}=\mathbf{1}\big[\,\mathrm{lower}(b(p))=\mathrm{lower}(b(g))\,\big]}
\begin{columns}[c,onlytextwidth]
\begin{column}{0.62\textwidth}
\worked{\texttt{Matrix.}\ $\neq$\ \texttt{Matrix?}\par\vspace{1mm}
 Yet $\mathrm{EM}=1$: $\nu$ strips the mark, $b$ keeps it.}
\end{column}
\begin{column}{0.34\textwidth}
\bigresult{0}{Break EM, punctuation only}
\end{column}
\end{columns}
\srcmono{break\_exact\_match\_rate $\cdot$ $b(x)$: spaces collapsed, \#k$\to$@@k@@, steps joined by @@SEP@@ $\cdot$ ADR 0026}
\end{frame}

% ------------------------------------------------------------ SARI (two slides)
\begin{frame}[t]
\slidehead{SETUP \textperiodcentered\ METRICS}{SARI: does it edit the question like gold?}
\planblock{QUESTION}{What languages are spoken in films \hlx{that} \hlx{share} \hlx{a} \hlx{director} \hlx{with} \textbf{The Matrix}?}
\planblock{GOLD PLAN}{\pline{1}{\hlt{Who} \hlt{directed} \textbf{The Matrix}?}%
 \pline{2}{\hlt{Which} \hlt{other} films \hlt{did} \hlt{[\#1]} \hlt{direct?}}%
 \pline{3}{What languages are spoken in \hlt{[\#2]?}}}
\vspace{-1mm}
\bigformula{\mathrm{SARI}=\tfrac13\big(\,\overline{\mathrm{keep}}+\overline{\mathrm{add}}+\overline{\mathrm{del}}\,\big)}
\legend{plain: kept\qquad \hlt{green}: added\qquad \hlx{red}: deleted\qquad $\bar{x}$: mean over $n$-grams, $n=1,\dots,4$}
\whyline{Rewards rewriting the question the way gold does.}
\srcmono{sari\_macro $\cdot$ split on spaces, case and punctuation kept $\cdot$ \#k$\to$@@k@@ $\cdot$ separator @@SEP@@ is an added word}
\end{frame}

\begin{frame}[t]
\slidehead{SETUP \textperiodcentered\ METRICS}{SARI, worked on the reversed plan}
\vspace{3mm}
\begin{columns}[c,onlytextwidth]
\begin{column}{0.58\textwidth}
\colhead{WORKED $\cdot$ REVERSED}\par\vspace{1mm}
\footnotesize
\begin{tabular}{@{}p{10mm}R{16mm}R{20mm}R{17mm}@{}}
\colhead{$n$} & \colhead{KEEP} & \colhead{ADD} & \colhead{DELETE}\\
\midrule
1 & \textcolor{inktwo}{8/8} & \textcolor{inktwo}{9/9} & \textcolor{inktwo}{5/5}\\
2 & \textcolor{inktwo}{5/5} & \textcolor{inktwo}{10/12} & \textcolor{inktwo}{7/7}\\
3 & \textcolor{inktwo}{3/3} & \textcolor{inktwo}{9/13} & \textcolor{inktwo}{8/8}\\
4 & \textcolor{inktwo}{2/2} & \textcolor{gold}{7/13} & \textcolor{inktwo}{8/8}\\
\midrule
mean & \textbf{1.00} & \textbf{0.766} & \textbf{1.00}\\
\bottomrule
\end{tabular}
\end{column}
\begin{column}{0.38\textwidth}
\bigresult{0.92}{$\tfrac13(1.00+0.766+1.00)$}
\end{column}
\end{columns}
\blindspot{Merged drops a step yet scores 0.77; an empty plan, 0.21.}
\srcmono{keep, add: F1; delete: precision; 0/0 = 1 $\cdot$ cell $t/s$ (here $s=r$) $\cdot$ reordering breaks cross-step $n$-grams}
\end{frame}

% ------------------------------------------------------------ GED (two slides)
\begin{frame}[t]
\slidehead{SETUP \textperiodcentered\ METRICS}{GED: how far is the plan's graph from gold's?}
\planpair{GOLD PLAN}{\goldlines}{EXTRA STEP ADDED}{\pline{1}{Who directed \textbf{The Matrix}?}%
 \pline{2}{Which other films did \plc{[\#1]} direct?}%
 \pline{3}{\hlx{What are the titles of those films?}}%
 \pline{4}{What languages are spoken in \plc{[\#2]}?}}
\vspace{1.5mm}
\tikzset{stepnode/.style={circle,draw=accent,line width=.5pt,minimum size=5.4mm,inner sep=0pt,
  font=\footnotesize\bfseries,text=accent},link/.style={-latex,inktwo,line width=.5pt}}
\begin{columns}[T,onlytextwidth]
\begin{column}{0.47\textwidth}
\begin{tikzpicture}[x=1mm,y=1mm]
  \node[stepnode] (g1) at (0,0) {1}; \node[stepnode] (g2) at (14,0) {2}; \node[stepnode] (g3) at (28,0) {3};
  \draw[link] (g2) -- (g1); \draw[link] (g3) -- (g2);
  \path (0,6) (0,-3);
\end{tikzpicture}
\end{column}
\begin{column}{0.47\textwidth}
\begin{tikzpicture}[x=1mm,y=1mm]
  \node[stepnode] (p1) at (0,0) {1}; \node[stepnode] (p2) at (14,0) {2};
  \node[stepnode,draw=bad,fill=badsoft,text=bad] (p3) at (28,0) {3}; \node[stepnode] (p4) at (42,0) {4};
  \draw[link] (p2) -- (p1); \draw[link] (p4) to[bend right=30] (p2);
  \path (0,6) (0,-3);
\end{tikzpicture}
\end{column}
\end{columns}
\vspace{-1mm}
\bigformula{\mathrm{GED}=\frac{\min_{\pi}\,c(\pi)}{\max\big(|V_p|+|E_p|,\ |V_g|+|E_g|\big)}}
\legend{node: a step\qquad edge $i\to k$: step $i$ cites \plc{\#k}\qquad $\pi$: cheapest edit path}
\whyline{One distance for wording, links and length. Lower is better.}
\srcmono{ged\_macro $\cdot$ a guard (node cap, time budget, no result) gives an upper bound: counted, not dropped}
\end{frame}

\begin{frame}[t]
\slidehead{SETUP \textperiodcentered\ METRICS}{GED, worked on the extra-step plan}
\vspace{3mm}
\begin{columns}[c,onlytextwidth]
\begin{column}{0.60\textwidth}
\colhead{WORKED $\cdot$ EXTRA STEP}\par\vspace{1mm}
\workedmath{
 \text{sizes} &{}= 5,\ 6 &\quad (3{+}2,\ 4{+}2)\\
 \text{steps } 1,2,4 &{}\mapsto 1,2,3 &\quad \text{cost } 0\\
 \text{delete step } 3 & &\quad \text{cost } 1\\
 \mathrm{GED} &{}= 1\,/\max(5,6)}
\end{column}
\begin{column}{0.36\textwidth}
\bigresult{0.17}{GED, extra step ($\downarrow$ better)}
\end{column}
\end{columns}
\blindspot{Reversed (0.40) beats merged (0.57): a graph has no step order.}
\srcmono{insert/delete 1 $\cdot$ edge relabel 0 $\cdot$ node relabel $1-2m/(|a|+|b|)$, $m$ aligned matching words $\cdot$ can exceed 1}
\end{frame}

% ------------------------------------------------------------ answer EM / F1 (one slide)
\begin{frame}[t]
\slidehead{SETUP \textperiodcentered\ METRICS}{Answer EM and F1: is the final answer right?}
\planpair{GOLD ANSWER}{English}{PREDICTED ANSWER}{\textcolor{muted}{the} English \hlx{language}}
\bigformula{\mathrm{EM}_A=\max_{a\in A}\mathbf{1}\big[\nu_A(\hat a)=\nu_A(a)\big]\qquad F_{1,A}=\max_{a\in A}\frac{2c}{|T(\hat a)|+|T(a)|}}
\legend{$T$: an answer's words\qquad $c$: words in common, as a multiset\qquad $A$: gold answer + aliases}
\begin{columns}[c,onlytextwidth]
\begin{column}{0.62\textwidth}
\worked{$\nu_A$: ``english language'' vs ``english'' $\Rightarrow$ $\mathrm{EM}_A=0$\par\vspace{1mm}
 $c=1$: $F_{1,A}=2\cdot1/(2+1)$}
\end{column}
\begin{column}{0.34\textwidth}
\bigresult{0\,/\,0.67}{answer EM / F1}
\end{column}
\end{columns}
\srcmono{src/answer\_metrics.py, MuSiQue official $\cdot$ $\nu_A$: lowercase, drop punctuation and a/an/the $\cdot$ ADR 0019}
\end{frame}

\begin{frame}[t]
\slidehead{SETUP \textperiodcentered\ METRICS}{The five plans on every plan metric}
\footnotesize
{\setlength{\tabcolsep}{2pt}%
\begin{tabular}{@{}p{20mm}*{3}{R{9.4mm}}R{11mm}*{7}{R{9.4mm}}@{}}
\colhead{PLAN} & \colhead{EM} & \colhead{STEP F1} & \colhead{ORD. ACC.} & \colhead{\mbox{ROUGE-L}}
 & \colhead{HOP EM} & \colhead{UNDER $\downarrow$} & \colhead{OVER $\downarrow$} & \colhead{CHAIN}
 & \colhead{BREAK EM} & \colhead{SARI} & \colhead{GED $\downarrow$}\\
\midrule
Reversed & \textcolor{inktwo}{0.00} & \textcolor{inktwo}{1.00} & \textcolor{inktwo}{0.33} & \textcolor{inktwo}{0.38} & \textcolor{inktwo}{1.00} & \textcolor{inktwo}{0.00} & \textcolor{inktwo}{0.00} & \textcolor{inktwo}{0.50} & \textcolor{inktwo}{0.00} & \textcolor{inktwo}{0.92} & \textcolor{inktwo}{0.40}\\
Merged & \textcolor{inktwo}{0.00} & \textcolor{inktwo}{0.40} & \textcolor{inktwo}{0.33} & \textcolor{inktwo}{0.64} & \textcolor{inktwo}{0.00} & \textcolor{inktwo}{1.00} & \textcolor{inktwo}{0.00} & \textcolor{inktwo}{1.00} & \textcolor{inktwo}{0.00} & \textcolor{inktwo}{0.77} & \textcolor{inktwo}{0.57}\\
Unlinked & \textcolor{inktwo}{0.00} & \textcolor{inktwo}{0.33} & \textcolor{inktwo}{0.33} & \textcolor{inktwo}{0.82} & \textcolor{inktwo}{1.00} & \textcolor{inktwo}{0.00} & \textcolor{inktwo}{0.00} & \textcolor{inktwo}{0.00} & \textcolor{inktwo}{0.00} & \textcolor{inktwo}{0.88} & \textcolor{inktwo}{0.49}\\
Extra step & \textcolor{inktwo}{0.00} & \textcolor{inktwo}{0.86} & \textcolor{inktwo}{0.50} & \textcolor{inktwo}{0.82} & \textcolor{inktwo}{0.00} & \textcolor{inktwo}{0.00} & \textcolor{inktwo}{1.00} & \textcolor{inktwo}{1.00} & \textcolor{inktwo}{0.00} & \textcolor{inktwo}{0.92} & \textcolor{inktwo}{0.17}\\
Punct.\ only & \textcolor{inktwo}{1.00} & \textcolor{inktwo}{1.00} & \textcolor{inktwo}{1.00} & \textcolor{inktwo}{0.94} & \textcolor{inktwo}{1.00} & \textcolor{inktwo}{0.00} & \textcolor{inktwo}{0.00} & \textcolor{inktwo}{1.00} & \textcolor{inktwo}{0.00} & \textcolor{inktwo}{0.90} & \textcolor{inktwo}{0.05}\\
\bottomrule
\end{tabular}}
\callout{Reversed ties for best on step F1 and is worst on ROUGE-L: no single metric suffices.}
\srcline{Every cell printed by Slides/score\_metric\_examples.py via the evaluator's score\_item $\cdot$ $\downarrow$ lower is better}
\end{frame}

% ============================================== 18 -- statistics ======
% ---- Every toy number is printed by Slides/score_metric_examples.py through the evaluator's
% ---- own _mcnemar / _paired_bootstrap / _paired_t_test_row.  Every REAL RUN number is read
% ---- verbatim from experiments/exp-005/metrics.json, comparisons.unguided_vs_oracle_guided
% ---- (Qwen3.5-9B, A = unguided, B = oracle-guided, the deck's hop-guidance slide).
\begin{frame}[t]
\slidehead{SETUP \textperiodcentered\ STATISTICS}{The metric's type chooses the test}
\begin{columns}[T,onlytextwidth]
\begin{column}{0.30\textwidth}
 \ptag{BINARY OUTCOMES}\phead{Exact McNemar}
 \pbody{EM, hop-count EM, Break EM, under, over.}
\end{column}
\begin{column}{0.30\textwidth}
 \ptag{CONTINUOUS SCORES}\phead{Paired bootstrap}
 \pbody{Step F1, ordered acc., ROUGE-L, SARI, GED, chain validity.}
\end{column}
\begin{column}{0.30\textwidth}
 \ptag{AGREEMENT CHECK}\phead{Paired t-test}
 \pbody{Every per-item metric, as a cross-check.}
\end{column}
\end{columns}
\callout{\textbf{Everything is paired.} Both arms answer the same questions, so per-question
 difficulty cancels out.}
\srcmono{ADR 0009, amended by 0017 $\cdot$ alpha 0.05 $\cdot$ evaluator: no correction; Holm afterwards where shown (exp-018)}
\end{frame}

\begin{frame}[t]
\slidehead{SETUP \textperiodcentered\ STATISTICS}{Exact McNemar: only the disagreements count}
\begin{columns}[T,onlytextwidth]
\begin{column}{0.56\textwidth}
\colhead{TOY $\cdot$ 10 QUESTIONS, \ck\ = CORRECT}\par\vspace{1mm}
\footnotesize
{\setlength{\tabcolsep}{2.2pt}\renewcommand{\arraystretch}{1.12}%
\begin{tabular}{@{}l*{10}{c}@{}}
 & 1 & 2 & 3 & 4 & 5 & 6 & 7 & 8 & 9 & 10\\
\midrule
A & \ck & \ck & \ck & \xx & \ck & \ck & \ck & \ck & \xx & \ck\\
B & \xx & \xx & \ck & \ck & \xx & \xx & \ck & \xx & \xx & \xx\\
 & \textcolor{accent}{\bfseries b} & \textcolor{accent}{\bfseries b} & & \textcolor{gold}{\bfseries c} & \textcolor{accent}{\bfseries b} & \textcolor{accent}{\bfseries b} & & \textcolor{accent}{\bfseries b} & & \textcolor{accent}{\bfseries b}\\
\end{tabular}}
\par\vspace{1mm}
{\centering\large$p=2\sum_{k=0}^{\min(b,c)}\dbinom{b+c}{k}\,2^{-(b+c)}$\par}
\end{column}
\begin{column}{0.40\textwidth}
\colhead{WORKED}\par\vspace{1mm}
\workedmath[1.3]{
 b,\ c &{}= 6,\ 1\\
 \textstyle\sum &{}= \tbinom70+\tbinom71 = 8\\
 p &{}= 2\cdot8\,/\,2^7}
\bigresult[24]{0.125}{not significant; underpowered}
\end{column}
\end{columns}
\calloutwarn{With 5 or fewer disagreements, $p\ge0.0625$: it can never reach 0.05.}
\srcreal{EM:\ $b=1$, $c=8$, uncorrected $p=0.039$}{$p$ capped at 1; $p=1$ if $b+c=0$ $\cdot$ underpowered if $N<30$ or min.\ attainable $p\ge\alpha$}
\end{frame}

\begin{frame}[t]
\slidehead{SETUP \textperiodcentered\ STATISTICS}{Paired bootstrap: resample the questions, keep the pairs}
\begin{columns}[T,onlytextwidth]
\begin{column}{0.56\textwidth}
\colhead{TOY $\cdot$ 5 QUESTIONS, STEP F1}\par\vspace{1mm}
\footnotesize
{\setlength{\tabcolsep}{3.2pt}%
\begin{tabular}{@{}l*{5}{c}c@{}}
 & 1 & 2 & 3 & 4 & 5 & mean\\
\midrule
A & 0.8 & 0.6 & 1.0 & 0.5 & 0.7 & \textbf{0.72}\\
B & 0.6 & 0.5 & 0.7 & 0.3 & 0.6 & \textbf{0.54}\\
\bottomrule
\end{tabular}}
\par\vspace{2.4mm}
\colhead{3 OF 10,000 RESAMPLES $\cdot$ SAME QUESTIONS FOR A AND B}\par\vspace{1mm}
{\setlength{\tabcolsep}{3.2pt}%
\begin{tabular}{@{}lr@{}}
questions 1 4 4 3 3 & $\Delta^*=0.76-0.52=0.24$\\
questions 5 1 4 2 1 & $\Delta^*=0.68-0.52=0.16$\\
questions 3 5 4 4 4 & $\Delta^*=0.64-0.44=0.20$\\
\end{tabular}}
\end{column}
\begin{column}{0.40\textwidth}
\colhead{CI = 2.5TH AND 97.5TH PERCENTILE OF $\Delta^*$}\par\vspace{1mm}
{\centering\large$\Delta=0.72-0.54=0.18$\par}
\vspace{2mm}
\bigresult[22]{[0.12,\ 0.24]}{$0\notin\mathrm{CI}$: significant}
\end{column}
\end{columns}
\srcreal{step F1:\ $\Delta=-0.024$, CI $[-0.035,-0.012]$}{$R$ = 10,000 resamples, seed 42 $\cdot$ $\Delta=A-B$ $\cdot$ the toy's $N=5<30$ would be flagged underpowered}
\end{frame}

\begin{frame}[t]
\slidehead{SETUP \textperiodcentered\ STATISTICS}{Paired t-test: a cross-check on the mean difference}
\begin{columns}[T,onlytextwidth]
\begin{column}{0.52\textwidth}
\colhead{SAME 5 QUESTIONS}\par\vspace{1mm}
\footnotesize
{\setlength{\tabcolsep}{3.2pt}%
\begin{tabular}{@{}l*{5}{c}@{}}
 & 1 & 2 & 3 & 4 & 5\\
\midrule
A & 0.8 & 0.6 & 1.0 & 0.5 & 0.7\\
B & 0.6 & 0.5 & 0.7 & 0.3 & 0.6\\
\midrule
$d=A-B$ & \textbf{0.2} & \textbf{0.1} & \textbf{0.3} & \textbf{0.2} & \textbf{0.1}\\
\bottomrule
\end{tabular}}
\par\vspace{3mm}
{\centering\large$\displaystyle t=\frac{\bar d}{s_d/\sqrt{N}},\qquad N-1\ \text{d.f.}$\par}
\end{column}
\begin{column}{0.44\textwidth}
\colhead{WORKED}\par\vspace{1mm}
\workedmath[1.3]{
 \bar d,\ s_d &{}= 0.18,\ 0.0837\\
 s_d/\sqrt5 &{}= 0.0374\\
 t &{}= 0.18\,/\,0.0374 = 4.81}
\vspace{1mm}
\bigresult[26]{0.009}{$p$, 4 d.f.: significant; underpowered}
\end{column}
\end{columns}
\srcreal{step F1:\ $t=-4.09$, uncorrected $p=4.8\times10^{-5}$}{scipy ttest\_rel $\cdot$ no claim if $N<2$ or $s_d=0$ $\cdot$ Holm (exp-018) applied later}
\end{frame}

% ======================================= 19 -- divider: results ======
{\setbeamercolor{background canvas}{bg=sunk}
\begin{frame}[t]
\vspace{18mm}
{\color{accent}\fontsize{11}{13}\selectfont Part three\par}
\vspace{4mm}
{\color{ink}\fontsize{24}{28}\selectfont Results\par}
\vspace{5mm}
{\color{inktwo}\fontsize{11}{15}\selectfont
 One claim per slide --- overall first, then broken out by hop depth.\par}
\end{frame}
}

% ========================================== 20 -- masking result ======
\begin{frame}[t]
\slidehead{RESULT \textperiodcentered\ MASKING}{Typed masking leads on every plan metric}
\small
\begin{tabular}{@{}p{26mm}R{20mm}R{20mm}R{20mm}R{20mm}R{20mm}@{}}
\colhead{VARIANT} & \colhead{EXACT MATCH} & \colhead{STEP F1} & \colhead{ORDERED ACC.}
 & \colhead{ROUGE-L F1} & \colhead{HOP-COUNT EM}\\
\midrule
typed & \textbf{0.058} & \textbf{0.201} & \textbf{0.179} & \textbf{0.544} & \textbf{0.522}\\
uniform & \textcolor{inktwo}{0.055} & \textcolor{inktwo}{0.190} & \textcolor{inktwo}{0.166} & \textcolor{inktwo}{0.543} & \textcolor{inktwo}{0.505}\\
raw & \textcolor{inktwo}{0.037} & \textcolor{inktwo}{0.178} & \textcolor{inktwo}{0.158} & \textcolor{inktwo}{0.532} & \textcolor{inktwo}{0.485}\\
\bottomrule
\end{tabular}
\calloutwarn{Margins over raw are consistent but \textbf{small} --- 2.1 points of exact match, 2.3 of step
 F1 --- and \textbf{no significance test has been run on typed vs uniform}. That ordering is
 provisional and is stated as provisional.}
\srcline{MuSiQue dev $\cdot$ n = 600 $\cdot$ unguided $\cdot$ Mistral-7B \hfill v1 prior work --- not re-run in v2}
\end{frame}

% ========================================= 21 -- per-hop decline ======
\begin{frame}[t]
\slidehead{RESULT \textperiodcentered\ BY HOP DEPTH}{Quality falls with every hop added}
\small
{\renewcommand{\arraystretch}{1.06}%
\begin{tabular}{@{}p{30mm}p{26mm}R{20mm}R{20mm}R{20mm}@{}}
\colhead{METRIC} & \colhead{VARIANT} & \colhead{2-HOP} & \colhead{3-HOP} & \colhead{4-HOP}\\
\midrule
Step F1 & typed & \textbf{0.322} & \textbf{0.171} & \textbf{0.109}\\
 & uniform & \textcolor{inktwo}{0.320} & \textcolor{inktwo}{0.157} & \textcolor{inktwo}{0.092}\\
 & raw & \textcolor{inktwo}{0.282} & \textcolor{inktwo}{0.147} & \textcolor{inktwo}{0.103}\\
\midrule
Hop-count EM & typed & \textbf{0.720} & \textbf{0.525} & \textcolor{inktwo}{0.320}\\
 & uniform & \textcolor{inktwo}{0.695} & \textcolor{inktwo}{0.495} & \textbf{0.325}\\
 & raw & \textcolor{inktwo}{0.680} & \textcolor{inktwo}{0.490} & \textcolor{inktwo}{0.285}\\
\bottomrule
\end{tabular}}
\callout{Step F1 loses roughly two-thirds of its value between 2 and 4 hops. Four-hop is both the
 hardest class and the one with the least training data --- those two facts are not
 independent. \textbf{Provenance:} this is the v1 measurement; the same typed arm re-run
 and re-scored in v2 sits about a point lower at every depth (0.710 / 0.515 / 0.300).}
\srcline{200 questions per hop $\cdot$ bold marks best in column \hfill v1 prior work --- not re-run in v2}
\end{frame}

% =============================================== 22 -- pool size ======
\begin{frame}[t]
\slidehead{RESULT \textperiodcentered\ POOL SIZE}{1,000 is too small; above that, nothing separates}
\vspace{2mm}
\begin{columns}[T,onlytextwidth]
\begin{column}{0.48\textwidth}
\begin{tikzpicture}[x=1mm,y=1mm]
  \barrow{24}{8,000}{100}{0.190}{accent}{38}
  \barrow{16}{4,000}{96.8}{0.184}{accent}{38}
  \barrow{8}{2,000}{95.8}{0.182}{accent}{38}
  \barrow{0}{1,000}{89.5}{0.170}{ruletwo}{38}
\end{tikzpicture}
\end{column}
\begin{column}{0.48\textwidth}
 \vspace{-2.6mm}
 \callout{Only the contrasts \textbf{involving 1,000} survive Holm correction. No contrast above
  1,000 survives on any metric --- at this sample size the adjacent steps are
  \textbf{underpowered, not equal}.}
 {\color{inktwo}\scriptsize Exact match stays flat across all cells regardless of size.\par}
\end{column}
\end{columns}
\srcpair{Mean step F1 $\cdot$ imbalanced arm}{v1 prior work, re-analysed 2026-08-20 --- not re-run in v2}
\end{frame}

% ==================================== 23 -- pool construction ======
\begin{frame}[t]
\slidehead{RESULT \textperiodcentered\ POOL CONSTRUCTION}{Construction changes everything upstream, and nothing downstream}
{\color{inktwo}\small
 Imbalanced against clustered, traced stage by stage. The two arms are physically separate
 the whole way --- and arrive at the same place.\par}
\vspace{4mm}
\begin{columns}[T,onlytextwidth]
\begin{column}{0.22\textwidth}
 \stat{accent}{0.053}{Jaccard overlap of the two \textbf{pools} --- 201 of 2,000 rows shared}
\end{column}
\begin{column}{0.23\textwidth}
 \stat{accent}{0 / 750}{queries whose \textbf{top-5 exemplars} are the same list. Most share none}
\end{column}
\begin{column}{0.22\textwidth}
 \stat{accent}{0 / 150}{logged \textbf{prompts} that match}
\end{column}
\begin{column}{0.23\textwidth}
 \stat{gold}{92\%}{of \textbf{generated plans} differ --- yet no reported metric moves detectably}
\end{column}
\end{columns}
\callout{\textbf{This was checked because it looked like a bug.} One cell had clustered and imbalanced
 tie at 0.5133 hop-count EM --- which turned out to be 258 discordant items splitting exactly
 129 / 129, in one cell of six. The other five do not tie. \textbf{It is a genuine null, not a
 duplicated run:} the same harness detects reranking's effect 9 times out of 9.}
\srcpair{exp-019 re-analysis of exp-010 $\cdot$ n = 750 per cell}{own dev sample, not the pinned 600 --- internally comparable only}
\end{frame}

% ================================================= 24 -- rerank ======
\begin{frame}[t]
\slidehead{RESULT \textperiodcentered\ RERANKING}{Reranking moves overlap and structure, not binary outcomes}
\vspace{1mm}
\begin{columns}[T,onlytextwidth]
\begin{column}{0.47\textwidth}
\begin{tikzpicture}[x=1mm,y=1mm]
  \barrow{49}{ROUGE-L}{100}{9/9 $\cdot$ 3\ck}{accent}{30}
  \barrow{42}{GED}{100}{9/9 $\cdot$ 4\ck}{accent}{30}
  \barrow{35}{SARI}{100}{9/9 $\cdot$ 0}{accent}{30}
  \barrow{28}{Step F1}{88.9}{8/9 $\cdot$ 1\ck}{accent}{30}
  \barrow{21}{Ordered}{77.8}{7/9 $\cdot$ 1\ck}{accent}{30}
  \barrow{14}{Hop-count}{88.9}{8/9 $\cdot$ 0}{ruletwo}{30}
  \barrow{7}{Break EM}{55.6}{5/9 $\cdot$ 0}{ruletwo}{30}
  \barrow{0}{Chain val.}{44.4}{4/9 $\cdot$ 0}{ruletwo}{30}
\end{tikzpicture}
\end{column}
\begin{column}{0.49\textwidth}
 \vspace{-2.6mm}
 \callout{\textbf{The earlier reading was wrong, not null.} ``Rerank doesn't move the score'' was
  measured on the retired composite, whose own intervals swing $\pm$0.2 on the very pairs
  where GED moves a stable +0.01--0.02. Wrong instrument, not absent effect.}
 {\color{inktwo}\scriptsize
  Pairs improved of 9, then Holm survivors. Pairs share one evaluation set, so sign
  counts are descriptive.\par}
\end{column}
\end{columns}
\srcline{exp-018 re-analysis of exp-010 $\cdot$ 9 matched pairs $\cdot$ n = 750 each}
\end{frame}

% ======================================== 25 -- oracle guidance ======
\begin{frame}[t]
\slidehead{RESULT \textperiodcentered\ HOP GUIDANCE}{A perfect hop count helps one model far more than the other}
\small
{\renewcommand{\arraystretch}{1.00}%
\begin{tabular}{@{}p{28mm}R{16mm}R{20mm}R{18mm}S{42mm}@{}}
\colhead{CONDITION} & \colhead{STEP F1} & \colhead{HOP-COUNT EM} & \colhead{EXACT MATCH}
 & \colhead{PAIRED TESTS VS UNGUIDED}\\
\midrule
\multicolumn{5}{@{}l}{\textbf{Mistral-7B-Instruct-v0.3}}\\
\midrule
Unguided & \textcolor{inktwo}{0.204} & \textcolor{inktwo}{0.508} & \textcolor{inktwo}{0.062} & \textcolor{muted}{baseline}\\
Oracle-guided & \textcolor{inktwo}{0.206} & \textcolor{gold}{\bfseries 0.590} & \textcolor{inktwo}{0.057}
 & \textcolor{inktwo}{count, ROUGE-L and GED move; content does not}\\
\midrule
\multicolumn{5}{@{}l}{\textbf{Qwen3.5-9B}}\\
\midrule
Unguided & \textcolor{inktwo}{0.213} & \textcolor{inktwo}{0.520} & \textcolor{inktwo}{0.062} & \textcolor{muted}{baseline}\\
Oracle-guided & \textcolor{gold}{\bfseries 0.237} & \textcolor{gold}{\bfseries 0.873} & \textcolor{gold}{\bfseries 0.073}
 & \textcolor{inktwo}{six of nine report-card metrics move, content and count included}\\
\bottomrule
\end{tabular}}
\vspace{-3.5mm}
\calloutopen{Same conditions, same 600 questions, same decoding. \textbf{Effectiveness is
 model-dependent}, which is why the next experiment varies only the model.}
\srcpair{exp-004 / exp-005, tested in exp-014}{not significant on Qwen: SARI, chain validity, Break EM}
\end{frame}

% ================================= 26 -- hop-matched retrieval ======
\begin{frame}[t]
\slidehead{RESULT \textperiodcentered\ HOP-MATCHED RETRIEVAL}{Matching the pool to the question's depth moves count, not quality}
\small
\begin{tabular}{@{}p{40mm}R{16mm}R{20mm}S{46mm}@{}}
\colhead{METRIC} & \colhead{MIXED POOL} & \colhead{ORACLE HOP-MATCHED} & \colhead{RESULT}\\
\midrule
Step F1 (macro) & \textcolor{inktwo}{0.204} & \textcolor{inktwo}{0.193}
 & \textcolor{muted}{not significant $\cdot$ CI [$-$0.006, +0.027]}\\
Ordered step accuracy & \textcolor{inktwo}{0.180} & \textcolor{inktwo}{0.180}
 & \textcolor{muted}{not significant $\cdot$ p = 0.950}\\
Hop-count exact match & \textcolor{inktwo}{0.515} & \textcolor{gold}{\bfseries 0.600}
 & \textcolor{inktwo}{significant $\cdot$ McNemar p = 0.0006}\\
Exact match & \textcolor{inktwo}{0.062} & \textcolor{inktwo}{0.057}
 & \textcolor{muted}{not significant $\cdot$ p = 0.664}\\
\bottomrule
\end{tabular}
\calloutwarn{The gain concentrates at 2-hop and also \textbf{shrinks the candidate pool} --- so it may be
 measuring pool size rather than hop matching. \textbf{Not adopted on this evidence.}}
\srcline{exp-009 $\cdot$ pinned 600 $\cdot$ unguided $\cdot$ Mistral-7B}
\end{frame}

% ========================================== 27 -- per-hop oracle ======
\begin{frame}[t]
\slidehead{RESULT \textperiodcentered\ BY HOP DEPTH}{Where the pool is thinnest, hop-matching starts to hurt}
\small
\begin{tabular}{@{}p{34mm}p{32mm}R{18mm}R{18mm}R{18mm}@{}}
\colhead{METRIC} & \colhead{ARM} & \colhead{2-HOP} & \colhead{3-HOP} & \colhead{4-HOP}\\
\midrule
Step F1 & mixed & \textcolor{inktwo}{0.327} & \textbf{0.168} & \textbf{0.116}\\
 & oracle-matched & \textbf{0.355} & \textcolor{inktwo}{0.153} & \textcolor{bad}{0.071}\\
\midrule
Ordered accuracy & mixed & \textcolor{inktwo}{0.318} & \textbf{0.138} & \textbf{0.085}\\
 & oracle-matched & \textbf{0.345} & \textcolor{inktwo}{0.135} & \textcolor{bad}{0.060}\\
\midrule
Hop-count EM & mixed & \textcolor{inktwo}{0.710} & \textcolor{inktwo}{0.525} & \textcolor{inktwo}{0.310}\\
 & oracle-matched & \textcolor{gold}{\bfseries 0.845} & \textcolor{gold}{\bfseries 0.580} & \textcolor{gold}{\bfseries 0.375}\\
\bottomrule
\end{tabular}
\callout{The count keeps improving at every depth. \textbf{Step quality inverts at 4 hops}, where
 restricting the pool leaves only the 1,175-row bucket to draw from.}
\srcpair{n = 200 per hop, both arms}{no per-hop significance test run --- the paired test is on the pooled 600}
\end{frame}

% ============================================ 28 -- fine-tuning ======
\begin{frame}[t]
\slidehead{RESULT \textperiodcentered\ FINE-TUNING}{Fine-tuning beats prompting at its best configuration}
\vspace{2mm}
\begin{columns}[T,onlytextwidth]
\begin{column}{0.50\textwidth}
\small
\begin{tabular}{@{}>{\raggedright\arraybackslash}p{34mm}R{14mm}R{17mm}@{}}
\colhead{ARM} & \colhead{STEP F1} & \colhead{ROUGE-L F1}\\
\midrule
Prompting, best pool & \textcolor{inktwo}{0.189} & \textcolor{inktwo}{0.545}\\
Fine-tuned LoRA & \textbf{0.341} & \textbf{0.665}\\
\bottomrule
\end{tabular}
\end{column}
\begin{column}{0.44\textwidth}
 \vspace{-3mm}
 \stat{accent}{10 / 10}{paired statistics significant --- seven bootstrap, three McNemar --- and
  every one of them favours the fine-tuned arm}
\end{column}
\end{columns}
\callout{The prompting arm here is prompting at its \textbf{strongest measured configuration}, not a
 strawman. Fine-tuning is the current leader, and the prompting study becomes the ablation
 behind it --- which is a finding, not a defeat.}
\srcline{exp-015 $\cdot$ pinned 600 $\cdot$ paired bootstrap + McNemar}
\end{frame}

% ========================================= 29 -- generalisation ======
\begin{frame}[t]
\slidehead{RESULT \textperiodcentered\ GENERALISATION}{The fine-tune transfers form, but not depth}
\vspace{3mm}
\begin{columns}[T,onlytextwidth]
\begin{column}{0.30\textwidth}
 \stat{accent}{200/200}{chains still well-formed at a depth never seen in training
  (chain validity 1.000)}
\end{column}
\begin{column}{0.30\textwidth}
 \stat{gold}{169/200}{but systematically \textbf{too short} --- under-decomposition 0.845}
\end{column}
\begin{column}{0.31\textwidth}
 \stat{bad}{14.5\%}{hop-count exact match on 4-hop, against \textbf{30.0\%} for the prompting
  baseline}
\end{column}
\end{columns}
\callout{Trained on 2- and 3-hop only, the model keeps producing \textbf{valid} chains at 4 hops ---
 just too few steps. Step F1 still favours the fine-tune (0.198 vs 0.117); hop-count
 flips to prompting. \textbf{Fine-tuning is dataset-bound in depth, not in form.}}
\srcline{exp-013 $\cdot$ LoRA trained on 2/3-hop, evaluated on the 4-hop slice $\cdot$ n = 200}
\end{frame}

% ============================================= 30 -- end to end ======
\begin{frame}[t]
\slidehead{RESULT \textperiodcentered\ END TO END}{Better plans do produce better answers --- up to three hops}
\small
\begin{tabular}{@{}p{48mm}R{18mm}R{18mm}R{38mm}@{}}
\colhead{DECOMPOSITION SOURCE} & \colhead{ANSWER EM} & \colhead{ANSWER F1}
 & \colhead{EM BY HOP (2 / 3 / 4)}\\
\midrule
None --- question asked whole & \textcolor{inktwo}{18.0\%} & \textcolor{inktwo}{27.2\%} & \textcolor{inktwo}{21.5 / 15.5 / 17.0}\\
Prompted, unguided & \textcolor{inktwo}{21.3\%} & \textcolor{inktwo}{29.5\%} & \textcolor{inktwo}{29.0 / 20.0 / 15.0}\\
Fine-tuned & \textcolor{inktwo}{24.2\%} & \textcolor{inktwo}{32.8\%} & \textcolor{inktwo}{34.0 / 24.0 / 14.5}\\
Gold decompositions (ceiling) & \textcolor{gold}{\bfseries 28.0\%} & \textcolor{gold}{\bfseries 37.3\%}
 & \textcolor{inktwo}{40.0 / 30.5 / {\color{bad}13.5}}\\
\bottomrule
\end{tabular}
\calloutopen{\textbf{Read the last column.} Answer accuracy rises monotonically overall --- but at 4 hops
 the ordering collapses entirely: all four cells sit between 13.5\% and 17.0\%, and asking
 the question whole is nominally \emph{best}. Decomposition pays off at 2--3 hops only.
 No significance tests between these cells yet.}
\srcline{exp-017 $\cdot$ same reader and context in every cell $\cdot$ 200 questions per hop}
\end{frame}

% =================================== 31 -- divider: where next ======
{\setbeamercolor{background canvas}{bg=sunk}
\begin{frame}[t]
\vspace{18mm}
{\color{accent}\fontsize{11}{13}\selectfont Part four\par}
\vspace{4mm}
{\color{ink}\fontsize{24}{28}\selectfont Where this goes next\par}
\vspace{5mm}
{\color{inktwo}\fontsize{11}{15}\selectfont
 What the evidence does not yet cover, and the experiment designed to cover it.\par}
\end{frame}
}

% ======================================= 32 -- the model confound ======
\begin{frame}[t]
\slidehead{NEXT \textperiodcentered\ THE GAP IN THE EVIDENCE}{Almost every result here is one model's result}
\begin{columns}[T,onlytextwidth]
\begin{column}{0.50\textwidth}
 {\color{inktwo}\small
  Of every MuSiQue arm run, all but one use \textbf{Mistral-7B}. The single Qwen arm
  disagrees with the Mistral finding it is usually quoted beside.\par}
 \calloutopen{So ``hop guidance is worth little'' and ``fine-tuning beats prompting'' are currently
  claims about \textbf{the same one model} --- and the model that disagrees with the first
  was never brought into the second.}
\end{column}
\begin{column}{0.45\textwidth}
 \ptag{THE EXPERIMENT}\phead{Hold everything, vary one thing}
 \pbody{Same 2,000-example imbalanced pool, same pinned 600, same unguided condition and
  decoding.}
 \vspace{2.5mm}

 \pillon{Mistral-7B}\hspace{1.5mm}\pillon{Qwen2.5-7B}\hspace{1.5mm}\pillon{Qwen2.5-14B}
 \vspace{1.5mm}

 {\color{inktwo}\scriptsize
  Mistral vs Qwen-7B isolates the \textbf{family}, Qwen-7B vs Qwen-14B the \textbf{size}.
  Each contrast moves exactly one factor.\par
  \vspace{1mm}
  Qwen3.5-9B is deliberately \emph{not} the size rung: it differs from Qwen2.5-7B in
  generation as well as parameter count, so it would confound the two. Bringing it into
  the fine-tuning claim is a separate, cheap arm.\par}
\end{column}
\end{columns}
\srcmono{open issue \#49 $\cdot$ \#50}
\end{frame}

% ============================================== 33 -- also next ======
\begin{frame}[t]
\slidehead{NEXT \textperiodcentered\ ALSO QUEUED}{Four more things, in order}
\begin{tabular}{@{}p{6mm}@{\hspace{2mm}}p{126mm}@{}}
\chainnum{1} & \textcolor{ink}{\fontsize{9.5}{12}\selectfont
 \textbf{A fine-tune trained on the pool alone} --- 2,000 examples, not the full 19,938 ---
 to isolate what the pool itself is worth.}\\[3mm]
\chainnum{2} & \textcolor{ink}{\fontsize{9.5}{12}\selectfont
 \textbf{A fixed 45-question anchor set}, 15 per hop, inspected by hand in every experiment
 from here on. The same questions every time.}\\[3mm]
\chainnum{3} & \textcolor{ink}{\fontsize{9.5}{12}\selectfont
 \textbf{Few-shot example ordering} --- currently highest-similarity-first, never varied.
 Reversed, and least-complex-first.}\\[3mm]
\chainnum{4} & \textcolor{ink}{\fontsize{9.5}{12}\selectfont
 \textbf{Whether the null generalises} --- pool construction is now measured as a genuine
 null, but masking, retriever and pool size have not had the same stage-by-stage check.}
\end{tabular}
\srcline{Every result from here on is reported overall \emph{and} per hop depth.}
\end{frame}

% ============================================ 34 -- limitations ======
\begin{frame}[t]
\slideheadplain{HONEST BOUNDARIES}{What this evidence does not support}
\begin{columns}[T,onlytextwidth]
\begin{column}{0.47\textwidth}
 \phead{Single-model scope}
 \pbody{The hop-guidance finding is Mistral-specific; Qwen3.5-9B behaves differently.}
 \vspace{3.5mm}
 \phead{Typed vs uniform untested}
 \pbody{The gap sits inside plausible run-to-run variance and carries no significance test.}
 \vspace{3.5mm}
 \phead{Pool evidence is one trial}
 \pbody{One draw per strategy mixes the strategy's effect with the noise of that draw.}
\end{column}
\begin{column}{0.47\textwidth}
 \phead{No human judgement}
 \pbody{Exact match against one gold plan penalises valid alternatives. No human validity
  study has been run.}
 \vspace{3.5mm}
 \phead{MetaQA end-to-end blocked}
 \pbody{The integrated graph run depends on an external system not yet available.}
 \vspace{3.5mm}
 \phead{No LLM judge, by design}
 \pbody{Early scoring used one; it is not ground truth and those figures are not evidence here.}
\end{column}
\end{columns}
\srcline{Anything not backed by a run artifact is reported as unmeasured.}
\end{frame}

% ============================================== 35 -- decisions ======
\begin{frame}[t]
\slidehead{TO AGREE ON}{Three decisions I need from this meeting}
\vspace{1mm}
\begin{tabular}{@{}p{6mm}@{\hspace{2mm}}p{126mm}@{}}
\chainnum{1} & \textcolor{ink}{\fontsize{9.5}{12}\selectfont
 \textbf{Is the report card the thesis's primary metric?}
 {\color{muted}Proposed, with end-to-end answer accuracy as the downstream complement.
 Brought to confirm, not to open.}}\\[3.4mm]
\chainnum{2} & \textcolor{ink}{\fontsize{9.5}{12}\selectfont
 \textbf{Is pool construction a contribution, or a hyperparameter?}
 {\color{muted}Proposal: a hyperparameter, fixed at size-2,000 imbalanced. It is now a
 measured null rather than an untested one --- reportable as an ablation, but it cannot be
 claimed as a gain.}}\\[3.4mm]
\chainnum{3} & \textcolor{ink}{\fontsize{9.5}{12}\selectfont
 \textbf{Do both datasets stay in scope, or does MuSiQue alone carry the thesis?}
 {\color{muted}MetaQA's end-to-end half is externally blocked.}}
\end{tabular}
\srcline{Everything already run stays in the thesis as the ablation behind these choices.}
\end{frame}

\end{document}
